% =====================================================================
%  CiteCheck: An Open Benchmark and Agentic Method for Verifiable
%             US Case-Law Citations
%
%  Target venue: NLLP @ EMNLP 2027 (primary) / TrustNLP @ EMNLP 2027
%
%  Requires the ACL 2024 style files (acl.sty, acl_natbib.bst). Download:
%      https://github.com/acl-org/acl-style-files
%  Place acl.sty and acl_natbib.bst in this directory (or have them on
%  Overleaf's class path) before building.
%
%  Build:        latexmk -pdf main.tex
%  Watch mode:   latexmk -pvc main.tex
%  Word count:   texcount main.tex
%
%  Conventions used in this skeleton:
%    \todo{...}  -- writing gaps (red text via \newcommand below; swap
%                   to the todonotes package if margin notes are wanted).
%    <TBD>       -- placeholder for a numeric experimental result.
%                   These are plain text so they grep cleanly.
%
%  This is a skeleton: prose is real where the outline supplied it;
%  numbers are <TBD> until Phase 3 experiments land.
% =====================================================================

\documentclass[11pt]{article}

% ----- ACL style (download from github.com/acl-org/acl-style-files) -----
% \usepackage[review]{acl}   % use during anonymous review
\usepackage{acl}             % camera-ready

% ----- Encoding and fonts -----
\usepackage[T1]{fontenc}
\usepackage[utf8]{inputenc}
\usepackage{microtype}

% ----- Math, tables, graphics -----
\usepackage{amsmath}
\usepackage{amssymb}
\usepackage{booktabs}
\usepackage{multirow}
\usepackage{graphicx}
\usepackage{url}

% ----- Hyperlinks (ACL style loads hyperref already; this guards
% against a second load on Overleaf if the user's template differs) -----
% \usepackage{hyperref}

% ----- Lightweight \todo{} so we don't need the todonotes package.
% Swap for \usepackage{todonotes} + \newcommand{\todo}[1]{\todonotes{#1}}
% if margin notes are preferred. -----
\usepackage{color}
\newcommand{\todo}[1]{\textcolor{red}{\textbf{[TODO:} #1\textbf{]}}}

% ----- Title -----
\title{CiteCheck: An Open Benchmark and Agentic Method for
       Verifiable US Case-Law Citations}

% ----- Authors (placeholder; ACL style uses \author + \And) -----
\author{John Paul L. Gabule \\
  \texttt{johnpaullimgabule@gmail.com} \\
  \And
  Advisor TBD \\
  \texttt{TBD@institution.edu} \\}

\begin{document}
\maketitle

% =====================================================================
%  Abstract
% =====================================================================
\begin{abstract}
% Skeleton beats: (1) hook, (2) empirical gap, (3) method gap,
% (4) contributions, (5) headline number, (6) artifact release.
In June 2023, two New York attorneys were sanctioned after filing a
brief containing six entirely fabricated case citations generated by a
large language model; the Charlotin AI-hallucination tracker has since
logged on the order of fifty US court orders flagging similar
incidents. The most rigorous published audit to date
\citep{magesh2025hallucination} measured hallucination rates of
17--33\% on leading commercial legal-research tools but did not
release the benchmark, prompts, or method, so the open community
cannot verify or improve on those numbers. Existing open legal
retrieval-augmented generation (RAG) work either targets non-US
jurisdictions \citep{wankhade2026elrag}, contracts rather than case
law \citep{pipitone2024legalbenchrag,butler2025mleb}, or verifies
citations post-hoc rather than at generation time
\citep{noel2025hallugraph,choi2025citeguard}. We introduce
\textbf{CiteCheck}, comprising: (a) an open benchmark of
${\sim}500$ US-grounded legal questions paired with gold citations and
per-citation labels on three orthogonal axes -- existence, support,
and jurisdictional validity; (b) a three-stage agentic-RAG pipeline
combining hybrid retrieval, a cross-encoder reranker trained with a
Bluebook-structure-aware groundedness loss parsed via \texttt{eyecite},
and an agent verify loop that resolves emitted citations against
CourtListener at decode time; and (c) an empirical comparison against
five open baselines, one closed reference, and one post-hoc detector,
plus four ablations. CiteCheck reduces fabrication rate from
\textbf{<TBD\_baseline>\%} to \textbf{<TBD\_ours>\%} on a single 24~GB
consumer GPU while maintaining \textbf{<TBD>} Answer Utility (5-point
Likert). We release the benchmark under CC-BY, code under MIT, and all
prompts and seeds at \url{https://github.com/<TBD>}.
\end{abstract}

% =====================================================================
%  1. Introduction
% =====================================================================
\section{Introduction}
\label{sec:intro}

\subsection{Opening: \emph{Mata v.\ Avianca}}
\label{sec:intro:mata}

In June 2023, Judge P.~Kevin Castel of the Southern District of New
York sanctioned two attorneys after they filed a brief containing six
entirely fabricated case citations -- ``Varghese v.\ China Southern
Airlines,'' ``Shaboon v.\ EgyptAir,'' and four others -- generated by
ChatGPT and presented as binding precedent
\citep{mata2023}. The opinions did not exist. Neither did the
reporter volumes that purportedly contained them. The episode made
plain that a hallucinated case citation is not an abstract NLP failure
mode but a sanctionable professional-responsibility failure that has
measurably injured real litigants and real lawyers.

\subsection{Scale of the problem}
\label{sec:intro:scale}

\emph{Mata} was not isolated. The Damien Charlotin AI-hallucination
tracker \citep{charlotin2025tracker} has logged on the order of fifty
US court orders since 2023 that flag AI-generated fabricated
citations, and the most rigorous empirical audit to date
\citep{magesh2025hallucination} found that even commercial
legal-research products marketed as ``hallucination-free'' --
Lexis+ AI, Thomson Reuters Westlaw AI-Assisted Research, and Practical
Law's Ask -- produce incorrect or unsupported information on 17\% to
33\% of legal queries. Comparable comparative audits
\citep{munir2025evaluating} confirm the order of magnitude.

\subsection{Why open and reproducible matters}
\label{sec:intro:open}

The populations most exposed to fabrication harm -- junior associates
billed out at first-year rates, public defenders, pro-bono and
\emph{pro se} litigants -- are also the populations least able to
afford the proprietary tools \citet{magesh2025hallucination} audited,
and least able to retain a senior partner to vet every cited
authority. \citet{cheong2025publicdefenders} document, through
interviews with US public defenders, that AI adoption in legal
services is gated less by appetite than by cost, auditability, and the
asymmetric reputational risk of a single bad citation. An open,
inspectable benchmark and method are preconditions for the populations
most exposed to fabrication harm to benefit from progress in
legal NLP at all.

\subsection{What is missing in existing work}
\label{sec:intro:gap}

Four gaps motivate this paper.

\paragraph{Audits are closed.}
\citet{magesh2025hallucination} measured hallucination rates on closed
commercial systems but released neither the benchmark, the prompts,
the baselines, nor the method. The open community cannot reproduce or
extend the audit.

\paragraph{Open legal RAG is non-US.}
EL-RAG \citep{wankhade2026elrag} is the closest open legal RAG with an
Evidence Alignment Layer, but it evaluates on
COLIEE/NyayaRAG/CaseLawBench, none of which is US case-law focused;
the paper is closed-access and there is no public code at the time of
writing.

\paragraph{US-focused open benchmarks miss citation faithfulness.}
LegalBench \citep{guha2022legalbench} and LegalBench-RAG
\citep{pipitone2024legalbenchrag} are open and US-focused but do not
isolate citation faithfulness as a scored axis; MLEB
\citep{butler2025mleb} evaluates embedding retrieval across six
jurisdictions but stops short of citation correctness.

\paragraph{Verification is post-hoc.}
HalluGraph \citep{noel2025hallugraph} and CiteGuard
\citep{choi2025citeguard} verify citations after the fact against a
knowledge graph or retrieval set; neither queries a live legal-opinion
registry inside the generation loop.

\subsection{Contributions}
\label{sec:intro:contrib}

We contribute three artifacts.

\begin{enumerate}
  \item \textbf{Benchmark.} CiteCheck v0.1, ${\sim}500$ US-grounded
        legal questions paired with gold citations and per-citation
        labels for existence, support, and jurisdictional validity.
        Released CC-BY with prompts, splits, and adversarial seed
        provenance drawn from \citet{charlotin2025tracker}.
  \item \textbf{Method.} A three-stage agentic-RAG pipeline whose two
        novel components are (a)~a cross-encoder reranker trained with
        a Bluebook-structure-aware groundedness loss,
        $\mathcal{L} = \mathcal{L}_{\text{rel}} + \lambda \cdot
        \mathcal{L}_{\text{ground}}$, parsed via \texttt{eyecite}
        \citep{cushman2018eyecite}; and (b)~an agent verify loop that
        emits citations through a constrained-decoding Bluebook
        grammar \citep{willard2023outlines} and resolves them at
        decode time against CourtListener with a local Caselaw Access
        Project mirror as fallback.
  \item \textbf{Empirical study.} Comparison against five open
        baselines (Vanilla 8B, Naive RAG, Self-RAG
        \citep{asai2023selfrag}, CRAG \citep{yan2024crag}, EL-RAG
        reimplementation), one closed reference (GPT-4o-mini with web
        search), and one post-hoc detector (HalluGraph), plus four
        ablations isolating each pipeline stage.
\end{enumerate}

\subsection{Roadmap}
\label{sec:intro:roadmap}

Section~\ref{sec:related} surveys the four bodies of work this paper
draws on. Section~\ref{sec:benchmark} specifies the benchmark.
Section~\ref{sec:method} specifies the method.
Sections~\ref{sec:setup} and~\ref{sec:results} report experimental
setup and results. Sections~\ref{sec:discussion} discusses
limitations, ethics, and future work; Section~\ref{sec:conclusion}
concludes.

% =====================================================================
%  2. Related Work
% =====================================================================
\section{Related Work}
\label{sec:related}

\subsection{Legal NLP benchmarks}
\label{sec:related:benchmarks}

Open evaluation of legal language models has, since 2021, coalesced
around a small number of benchmarks: CUAD
\citep{hendrycks2021cuad} for contract clause classification,
LegalBench \citep{guha2022legalbench} for multi-task legal reasoning,
LegalBench-RAG \citep{pipitone2024legalbenchrag} for retrieval-
augmented contract QA, MLEB \citep{butler2025mleb} for legal embedding
evaluation across six jurisdictions, the follow-up Legal RAG Bench
\citep{butler2026legalragbench} for end-to-end legal RAG, and the
recent CLAUSE benchmark \citep{choudhury2025clause} for auditing legal
reasoning capabilities; smaller-model legal reasoning has been studied
in \citep{vaddi2026smallmodels}. None of these benchmarks isolates
``is this citation real, supportive, and jurisdictionally
appropriate'' as a first-class evaluation axis, and most concentrate
on contracts rather than case-law-grounded research questions.
LegalBench-RAG's ``answer faithfulness'' metric is closer in spirit
than the others but is still scored over passages, not over emitted
citations.

\subsection{Legal RAG systems}
\label{sec:related:legalrag}

Recent open legal RAG systems improve specific pipeline stages.
\citet{reuter2025reliableretrieval} propose Summary-Augmented Chunking
to reduce document-level retrieval mismatch on large legal corpora;
\citet{topollai2025contracts} build a modular industrial
contract-management pipeline; \citet{wankhade2026elrag} introduces
EL-RAG with hybrid sparse-dense retrieval, learning-to-rank, multi-hop
reasoning, and an Evidence Alignment Layer. EL-RAG is the closest
published comparator to CiteCheck on legal multi-hop reasoning, but it
does not address citation existence and is evaluated on non-US
benchmarks. \citet{jia2025legalkg} and
\citet{dechtiar2025graphgrpolex} introduce legal-knowledge-graph and
contract-graph-plus-RL variants respectively; neither targets US case-
law citation existence. \citet{chen2025l4m} pursue a similar
trust-through-formal-reasoning direction.

\subsection{Citation faithfulness in general LLMs}
\label{sec:related:faithfulness}

Outside law, ALCE \citep{gao2023alce} established the modern
formulation of citation-aware generation by treating citation
correctness as a scored output property. Subsequent work --
\citet{ji2024chainofthought} on chain-of-thought citation,
\citet{wu2025sourcecheckup} on SourceCheckup for medical references,
\citet{choi2025citeguard} on CiteGuard for scientific writing,
\citet{tan2025interpdetect} on interpretable RAG-hallucination signals,
and \citet{su2025unifact} on hallucination/fact-verification
unification -- verifies that citations are well-supported by retrieved
evidence. None addresses the case-law-specific failure modes of
citation \emph{existence} (the cited reporter volume does not exist)
and \emph{jurisdictional validity} (the cited opinion is from a
non-binding court). \citet{guan2025faithfulnessaware} is the closest
methodological prior for our Bluebook-aware rerank loss, though that
work is domain-general; \citet{khan2026reasonverify} pursues a related
reason-and-verify framing.

\subsection{Agentic RAG architectures}
\label{sec:related:agentic}

Self-RAG \citep{asai2023selfrag} introduced self-reflective control
tokens for retrieve/no-retrieve and support-level decisions; CRAG
\citep{yan2024crag} added a retrieval evaluator that triggers
corrective re-retrieval; Adaptive-RAG \citep{jeong2024adaptiverag}
adapts retrieval depth to estimated question complexity; Self-CRAG
\citep{anand2026selfcrag} extends the corrective pattern to
conversational settings. More recent planner-centric work -- Beyond
ReAct \citep{wei2025beyondreact}, ReCAP \citep{zhang2025recap}, and
Lateral Tree-of-Thoughts \citep{madahar2025lateraltot} -- improves the
planning loop on general-domain agent benchmarks. None of these
architectures has been instantiated with a legal-specific verification
tool such as a live opinion-registry lookup.

\subsection{Hallucination detection in legal RAG}
\label{sec:related:legalhallucination}

The two empirical works most directly adjacent to CiteCheck are
\citet{magesh2025hallucination}, which audits closed commercial tools,
and HalluGraph \citep{noel2025hallugraph}, which proposes auditable
hallucination detection for legal RAG via knowledge-graph alignment.
\citet{magesh2025hallucination} is closed: no released benchmark or
method. HalluGraph is open but operates as a post-hoc detector on top
of an arbitrary generator, not as a generation-time verifier.
\citet{munir2025evaluating} provides a complementary comparative audit
of legal-operations tools. \citet{jha2025kglegal} pushes the
knowledge-graph framing into specific statutes (Title~18 \S 175).
\citet{savelka2023unreasonable} document the zero-shot strength of
LLMs on legal-text annotation tasks but do not isolate citation
existence. CiteCheck differs from \citet{magesh2025hallucination} by
being fully open and reproducible, and from HalluGraph by closing the
verification loop at decode time rather than after generation.

\subsection{What CiteCheck contributes that prior work does not}
\label{sec:related:positioning}

Table~\ref{tab:positioning} positions CiteCheck against the closest
prior work along six properties: open data, open code, US case-law
focus, an explicit citation-existence axis, an explicit jurisdiction
axis, and decode-time (as opposed to post-hoc) verification.

\begin{table*}[t]
\centering
\small
\begin{tabular}{lcccccc}
\toprule
\textbf{Work} & \textbf{Open data} & \textbf{Open code} &
\textbf{US case-law} & \textbf{Existence axis} &
\textbf{Jurisdiction axis} & \textbf{Decode-time verify} \\
\midrule
LegalBench / LegalBench-RAG     & Y & Y & Y/Partial & N & N & N \\
MLEB                            & Y & Y & Partial   & N & N & N \\
EL-RAG                          & N & N & N         & N & N & N \\
\citet{magesh2025hallucination} & N & N & Y         & Partial & N & N \\
CiteGuard / SourceCheckup       & Y & Y & N (sci/med) & Partial & N & Partial \\
HalluGraph                      & Y & Y & Y         & Y (post-hoc) & N & N \\
\textbf{CiteCheck (this work)}  & \textbf{Y} & \textbf{Y} & \textbf{Y} &
                                  \textbf{Y} & \textbf{Y} & \textbf{Y} \\
\bottomrule
\end{tabular}
\caption{Positioning of CiteCheck against the closest open and closed
prior work. ``Decode-time verify'' means a citation registry is
queried inside the generation loop rather than after the answer is
returned.}
\label{tab:positioning}
\end{table*}

% =====================================================================
%  3. The CiteCheck Benchmark
% =====================================================================
\section{The CiteCheck Benchmark}
\label{sec:benchmark}

\subsection{The three axes}
\label{sec:benchmark:axes}

CiteCheck scores each emitted citation on three orthogonal axes.

\paragraph{Existence (binary).}
Does the cited authority (case name + reporter + court + year)
resolve to a real opinion in CourtListener \citep{courtlistener},
with the Caselaw Access Project \citep{cap2018} bulk mirror as
fallback? This axis is machine-checkable.

\paragraph{Support (three-valued, binarized for headline).}
Given the proposition that the answer asserts the citation stands for,
does the resolved opinion text actually entail or substantiate that
proposition? Labels are
$\{\text{entails}, \text{neutral}, \text{contradicts}\}$, binarized to
$\{\text{supports}, \text{does not support}\}$ for the primary metric.

\paragraph{Jurisdictional validity (three-valued).}
Given the forum and the legal question, is the cited court
\emph{binding}, \emph{persuasive}, or \emph{off-point}?

\todo{add one worked example per axis: a question + a candidate
citation that is positive on one axis and negative on another, to
make the three-axis decomposition concrete.}

\subsection{Construction methodology}
\label{sec:benchmark:construction}

\paragraph{Seed.}
Adversarial questions are seeded from the Charlotin AI-hallucination
tracker \citep{charlotin2025tracker}. For each logged sanction order,
we extract the underlying legal question and one fabricated citation
and convert the pair into a question paired with a known-negative
citation, giving us a corpus of attested failure modes rather than
synthetic ones.

\paragraph{Corpus side.}
The verification corpus is the Caselaw Access Project bulk download
\citep{cap2018} (${\sim}6.7$M opinions) for offline retrieval, with
the CourtListener API \citep{courtlistener} as the live registry of
record for verification and metadata.

\paragraph{LLM pre-label.}
Following the cost-efficient labeling methodology of
\citet{gray2024llmlabeling}, GPT-4o-mini drafts per-axis labels for
approximately 1,500 candidate question/citation pairs.

\paragraph{Human verify.}
We subset the pre-labeled candidates to ${\sim}500$ pairs for
human-verified gold and double-annotate 100 of them to estimate
inter-annotator agreement.

\subsection{Coverage}
\label{sec:benchmark:coverage}

CiteCheck v0.1 contains $N = <TBD>$ questions (target $\approx 500$),
distributed across court tiers as shown in Table~\ref{tab:coverage}.
The v0.1 release restricts to federal appellate (the thirteen US
Circuit Courts of Appeals) and the US Supreme Court. State appellate
coverage is deferred to v0.2 owing to uneven CourtListener coverage at
the state level. Topical coverage spans civil procedure, contracts
(with overlap to CUAD \citep{hendrycks2021cuad}), constitutional law,
and criminal procedure.

\begin{table}[t]
\centering
\small
\begin{tabular}{lrr}
\toprule
\textbf{Court tier} & \textbf{Items} & \textbf{\%} \\
\midrule
US Supreme Court            & <TBD> & <TBD> \\
Federal Circuits (1--11)    & <TBD> & <TBD> \\
DC Circuit                  & <TBD> & <TBD> \\
Federal Circuit             & <TBD> & <TBD> \\
\midrule
\textbf{Total (v0.1)}       & \textbf{<TBD>} & \textbf{100.0} \\
\bottomrule
\end{tabular}
\caption{Court-tier coverage of CiteCheck v0.1. State appellate is
deferred to v0.2.}
\label{tab:coverage}
\end{table}

\subsection{Annotation protocol and agreement}
\label{sec:benchmark:annotation}

The annotation interface, instructions, and 50-item gold-calibration
set are described in Appendix~\ref{app:annotation}. On the 100-item
double-annotated subset, inter-annotator agreement (Cohen's $\kappa$)
is: existence $\kappa = <TBD>$, support $\kappa = <TBD>$,
jurisdictional validity $\kappa = <TBD>$.

\subsection{Comparison to closest existing benchmarks}
\label{sec:benchmark:comparison}

Table~\ref{tab:benchmark-comparison} compares CiteCheck against the
closest open benchmarks on the features that matter for case-law
citation verification: per-citation rather than passage-level labels,
and explicit existence, support, and jurisdiction axes.

\begin{table*}[t]
\centering
\small
\begin{tabular}{lp{2.3cm}p{1.6cm}p{2.2cm}cccc}
\toprule
\textbf{Benchmark} & \textbf{Items} & \textbf{US case-law?} &
\textbf{Per-citation labels?} & \textbf{Exist.} &
\textbf{Supp.} & \textbf{Jur.} \\
\midrule
CUAD              & 13k clauses, 510 contracts & No (contracts)
                  & Span-level     & N & N & N \\
LegalBench        & ${\sim}162$ tasks          & Partial
                  & Task-level     & N & N & N \\
LegalBench-RAG    & ${\sim}7$k QA over 700 clauses & Contracts
                  & Passage-level  & N & Partial & N \\
MLEB              & 10 retrieval sets, 6 jur.  & Multi-juris.
                  & Retrieval-pair & N & N & N \\
HalluGraph eval   & <TBD>                      & Partial
                  & Per-claim      & Post-hoc & Partial & N \\
\textbf{CiteCheck v0.1} & \textbf{${\sim}500$ QA} & \textbf{Yes}
                  & \textbf{Per-citation, 3-axis}
                  & \textbf{Y} & \textbf{Y} & \textbf{Y} \\
\bottomrule
\end{tabular}
\caption{Comparison of CiteCheck v0.1 to closest open legal-NLP
benchmarks on case-law focus, label granularity, and axis coverage.}
\label{tab:benchmark-comparison}
\end{table*}

% =====================================================================
%  4. Method: The CiteCheck Pipeline
% =====================================================================
\section{Method: The CiteCheck Pipeline}
\label{sec:method}

\subsection{Architecture overview}
\label{sec:method:overview}

The CiteCheck pipeline has three stages, illustrated in
Figure~\ref{fig:architecture}: (1) hybrid sparse-plus-dense retrieval,
(2) a faithfulness-tuned cross-encoder reranker, and (3) an agent
verify loop that emits Bluebook-grammar-conformant citations and
resolves each one against a live opinion registry before returning the
final answer. The implementation lives in the open-source
\texttt{citecheck} package, organized as
\texttt{retrieval/}, \texttt{reranker/}, \texttt{agent/},
\texttt{baselines/}, \texttt{data/}, and \texttt{eval/} submodules.
We refer the reader to the released design specification
(\texttt{project/citecheck\_design.md}) for full hyperparameter and
training detail.

\begin{figure*}[t]
\centering
% Real figure goes in figures/architecture.pdf; see Phase 3 Task 42.
% \includegraphics[width=0.95\textwidth]{figures/architecture.pdf}
\fbox{\parbox[c][7em][c]{0.85\textwidth}{
  \centering
  \textbf{[Figure placeholder]} 3-stage CiteCheck architecture:
  question $\rightarrow$ Stage~1 hybrid retrieve (BM25 + BGE,
  RRF-fused) $\rightarrow$ Stage~2 faithfulness-tuned cross-encoder
  rerank $\rightarrow$ Stage~3 agent verify loop (constrained
  decoding, \texttt{CitationResolver}, retrieve-on-failure)
  $\rightarrow$ annotated answer.
}}
\caption{Three-stage CiteCheck pipeline. Citations emitted in
Stage~3 are tagged \texttt{verified}, \texttt{unresolvable}, or
\texttt{non\_supporting} before being returned to the user.}
\label{fig:architecture}
\end{figure*}
\todo{generate figure: figures/architecture.pdf (Phase 3 Task 42).}

\subsection{Stage 1: Hybrid retrieval}
\label{sec:method:retrieval}

\paragraph{Sparse.}
BM25 \citep{robertson2009bm25} indexing of the local CAP bulk
download via Pyserini.

\paragraph{Dense.}
A dual encoder; the primary backbone is
\texttt{BAAI/bge-base-en-v1.5} \citep{xiao2024cpack}, A/B compared
against \texttt{nlpaueb/legal-bert-base-uncased}
\citep{chalkidis2020legalbert}.

\paragraph{Fusion.}
Reciprocal Rank Fusion over top-50 candidates from each branch
returns top-20 to the reranker.

This follows the standard retrieve-then-rerank framing of
\citet{lewis2020rag}, with a closer rerank-then-link analogue in the
biomedical entity-linking work of \citet{cao2025retrievethenrerank}.

\subsection{Stage 2: Faithfulness-tuned reranker}
\label{sec:method:reranker}

The reranker is a cross-encoder fine-tuned with a two-term loss,
\begin{equation}
\mathcal{L} = \mathcal{L}_{\text{rel}}
            + \lambda \cdot \mathcal{L}_{\text{ground}},
\label{eq:loss}
\end{equation}
where $\mathcal{L}_{\text{rel}}$ is the standard pointwise relevance
objective and $\mathcal{L}_{\text{ground}}$ is a binary cross-entropy
on whether the candidate passage \emph{contains a Bluebook-parseable
citation whose resolved opinion text supports the asserted claim},
parsed via the \texttt{eyecite} Bluebook-citation parser
\citep{cushman2018eyecite}. We sweep $\lambda \in \{0.1, 0.3, 0.5,
1.0\}$ on the development split; the chosen value is
$\lambda = <TBD>$. Mining of $(\text{query}, \text{doc},
\text{label})$ triples produces 10k--30k triples from the CiteCheck
dev split with Stage-1 BM25 candidates as hard and easy negatives.
The closest methodological priors are
\citet{guan2025faithfulnessaware} (faithfulness-aware multi-objective
ranking), \citet{khan2026reasonverify} (reason-and-verify framing),
and \citet{li2023parade} (passage representation aggregation for
reranking).

To our reading of the literature, the use of the \emph{Bluebook
citation structure} itself -- parsed reporter, court, year, pinpoint
-- as a groundedness signal in the rerank loss is novel.

\subsection{Stage 3: Agent verify loop}
\label{sec:method:agent}

The Stage-3 agent has three sub-components.

\paragraph{Constrained decoding.}
A Bluebook grammar specified for the \texttt{outlines} library
\citep{willard2023outlines} covers full case cites, short cites,
\texttt{id.}, \texttt{supra}, and parallel cites. The v0.1 grammar
covers ${\sim}80\%$ of common US case-cite forms encountered in the
CAP development sample; the residual 20\% is logged and tabulated in
Appendix~\ref{app:grammar}. We adopt the latency/quality trade-off
documented by \citet{schall2025constraineddecoding} as guidance for
when to relax the constraint.

\paragraph{\texttt{CitationResolver} tool.}
A tool the agent invokes that queries the CourtListener API
\citep{courtlistener} and falls back to the local CAP mirror
\citep{cap2018} on rate-limit. It returns the tuple
$(\text{exists}, \text{opinion\_text}, \text{court}, \text{year})$
to the agent. A persistent cache layer collapses duplicate lookups
within a session.

\paragraph{Verify loop.}
Up to $k = 3$ iterations of
(retrieve $\rightarrow$ generate $\rightarrow$ resolve $\rightarrow$
entail-check $\rightarrow$ re-retrieve on failure). Citations are
tagged \texttt{verified}, \texttt{unresolvable}, or
\texttt{non\_supporting} in the final answer. The pattern follows
Self-RAG \citep{asai2023selfrag} and CRAG \citep{yan2024crag},
extended by \citet{anand2026selfcrag} to conversational settings;
ours differs in calling a domain-specific live registry rather than a
generic search engine. Independent semantic-verification work
\citep{martin2024semanticverification,yue2024fact} and clinical
agentic systems \citep{low2025clinical} provide methodological
parallels.

\subsection{Novelty discussion}
\label{sec:method:novelty}

\paragraph{Live-registry verification vs. post-hoc detection.}
\citet{magesh2025hallucination} and HalluGraph
\citep{noel2025hallugraph} operate on completed model outputs.
CiteCheck instead queries CourtListener inside the generation loop, so
unresolvable citations trigger a retraction-or-retry before the user
sees them. This converts citation faithfulness from a downstream
audit into a generation-time hard constraint.

\paragraph{Structure-aware reranking.}
Existing faithfulness rerankers
\citep{guan2025faithfulnessaware,khan2026reasonverify} score relevance
and faithfulness over free text. CiteCheck's reranker additionally
exploits the typed-field structure exposed by \texttt{eyecite} --
reporter, court, year, pinpoint -- so the loss can attend to whether
the \emph{correct kind} of authority appears in the candidate
passage, not merely a textually similar one. This complements
related ontology-grounded RAG work \citep{clarizia2025ontologyrag} and
paragraph-perspective case-law citation prediction
\citep{olsen2023paragraphcaselaw}.

% =====================================================================
%  5. Experimental Setup
% =====================================================================
\section{Experimental Setup}
\label{sec:setup}

\subsection{Baselines}
\label{sec:setup:baselines}

We compare CiteCheck against the seven systems listed in
Table~\ref{tab:baselines}. The protocol contract every baseline must
implement is fixed at the file-level (\texttt{baselines/protocol.py})
to prevent leakage of CiteCheck-specific assumptions into the
comparison.

\begin{table*}[t]
\centering
\small
\begin{tabular}{llll}
\toprule
\textbf{System} & \textbf{Type} & \textbf{Reference}
                & \textbf{Code path} \\
\midrule
Vanilla Llama-3.1-8B-Instruct & No retrieval
   & Meta AI 2024
   & \texttt{baselines/vanilla.py} \\
Naive RAG (BM25 + top-$k$ stuffing) & Static RAG
   & \citet{lewis2020rag}
   & \texttt{baselines/naive\_rag.py} \\
Self-RAG & Self-reflective agentic RAG
   & \citet{asai2023selfrag}
   & \texttt{baselines/self\_rag.py} \\
CRAG & Corrective RAG
   & \citet{yan2024crag}
   & \texttt{baselines/crag.py} \\
EL-RAG (reimpl.) & Legal multi-hop RAG
   & \citet{wankhade2026elrag}
   & \texttt{baselines/el\_rag.py} \\
GPT-4o-mini + web search & Closed-API reference
   & OpenAI
   & n/a (vendor API) \\
HalluGraph (post-hoc) & Post-hoc detector
   & \citet{noel2025hallugraph}
   & n/a (scorer over saved outputs) \\
\bottomrule
\end{tabular}
\caption{Baselines compared in this paper.
``post-hoc'' HalluGraph runs over the saved outputs of Naive RAG.}
\label{tab:baselines}
\end{table*}

The EL-RAG reimplementation is a documented reproduction risk
(closed-access paper, no public code); any gap relative to the
description in \citet{wankhade2026elrag} is reported in
Appendix~\ref{app:repro}. If reproduction fails the substitute
comparator is IRCoT.

\subsection{Generators}
\label{sec:setup:generators}

\paragraph{Primary backbone.}
Llama-3.1-8B-Instruct, 4-bit QLoRA where applicable, run on a single
24~GB consumer GPU.

\paragraph{Secondary backbone.}
Qwen2.5-7B-Instruct (4-bit QLoRA) to test that gains are not
Llama-specific.

\paragraph{Stretch backbone (compute permitting).}
Llama-3.1-13B-Instruct on Colab Pro A100 spot.

\subsection{Compute}
\label{sec:setup:compute}

All primary experiments are run on a single 24~GB consumer GPU, with
QLoRA fine-tuning via \texttt{bitsandbytes} and \texttt{peft}. Bursts
to Colab Pro A100 are used for the reranker $\lambda$-sweep and for
the 13B stretch. Total compute budget is $<TBD>$ GPU-hours.

\subsection{Evaluation set}
\label{sec:setup:evalset}

The primary evaluation set is CiteCheck v0.1 (${\sim}500$ questions
with per-axis gold labels; Section~\ref{sec:benchmark}). Two
human-audit subsets are drawn: a 200-item subset for Citation Support
F1 against the NLI judge, and a 100-item subset for Answer Utility
on a 5-point Likert scale. Following the agentic-evaluation
methodology of \citet{gao2025ragalyst}, NLI scores are reported
alongside, not in place of, human-audited scores.

\subsection{Random seeds and stability}
\label{sec:setup:seeds}

We use three random seeds: 13, 42, and 1337. For CiteCheck and the
two strongest baselines, all three seeds are run on the full
evaluation set. For the remaining baselines we report a single seed
($=42$) unless budget permits more; this asymmetry is flagged
transparently in the per-table notes.

\subsection{Statistical reporting}
\label{sec:setup:stats}

We report mean $\pm$ standard deviation across seeds for every cell
where two or more seeds were run. For headline metrics (Fabrication
Rate, Citation Support F1) we report a paired bootstrap (10k
resamples) at the item level for CiteCheck against each baseline.

% =====================================================================
%  6. Results
% =====================================================================
\section{Results}
\label{sec:results}

All numeric values reported in this section are placeholders; final
values will be drawn from \texttt{citecheck/experiments/} in Phase~3.

\subsection{Headline results}
\label{sec:results:headline}

Table~\ref{tab:headline} reports the five headline metrics across all
systems. As a contextual reference (\emph{not} a direct comparison,
since the benchmark differs), \citet{magesh2025hallucination} measured
hallucination rates of 17--33\% on Lexis+~AI, Westlaw, and Practical
Law on their proprietary benchmark.

\begin{table*}[t]
\centering
\small
\begin{tabular}{lccccc}
\toprule
\textbf{System} &
  \textbf{Fabrication $\downarrow$} &
  \textbf{Resolution $\uparrow$} &
  \textbf{Support F1 $\uparrow$} &
  \textbf{Jur. Validity $\uparrow$} &
  \textbf{Utility (Likert) $\uparrow$} \\
\midrule
Vanilla Llama-3.1-8B            & <TBD> & <TBD> & <TBD> & <TBD> & <TBD> \\
Naive RAG                       & <TBD> & <TBD> & <TBD> & <TBD> & <TBD> \\
Self-RAG                        & <TBD> & <TBD> & <TBD> & <TBD> & <TBD> \\
CRAG                            & <TBD> & <TBD> & <TBD> & <TBD> & <TBD> \\
EL-RAG (reimpl.)                & <TBD> & <TBD> & <TBD> & <TBD> & <TBD> \\
GPT-4o-mini + web               & <TBD> & <TBD> & <TBD> & <TBD> & <TBD> \\
\textbf{CiteCheck (ours)}       & \textbf{<TBD>} & \textbf{<TBD>}
                                & \textbf{<TBD>} & \textbf{<TBD>}
                                & \textbf{<TBD>} \\
\bottomrule
\end{tabular}
\caption{Headline results on CiteCheck v0.1, mean over 3 seeds where
applicable. Fabrication Rate counts emitted citations that fail the
Existence axis; Citation Resolution Rate is its complement on
emitted-citation count. Support~F1 is computed against the NLI judge;
Section~\ref{sec:results:audit} reports the human-audited variant on
the 200-item subset.}
\label{tab:headline}
\end{table*}

\subsection{Per-axis breakdown}
\label{sec:results:peraxis}

Table~\ref{tab:peraxis} decomposes performance by axis for CiteCheck
and the top three baselines.

\begin{table*}[t]
\centering
\small
\begin{tabular}{lcccc}
\toprule
\textbf{System} &
  \textbf{Existence Acc.} &
  \textbf{Support F1 (NLI)} &
  \textbf{Support F1 (human, $n{=}200$)} &
  \textbf{Jur. Validity} \\
\midrule
Top baseline \#1     & <TBD> & <TBD> & <TBD> & <TBD> \\
Top baseline \#2     & <TBD> & <TBD> & <TBD> & <TBD> \\
Top baseline \#3     & <TBD> & <TBD> & <TBD> & <TBD> \\
\textbf{CiteCheck}   & \textbf{<TBD>} & \textbf{<TBD>}
                     & \textbf{<TBD>} & \textbf{<TBD>} \\
\bottomrule
\end{tabular}
\caption{Per-axis scores for CiteCheck and the three strongest
baselines.}
\label{tab:peraxis}
\end{table*}

\subsection{Per-jurisdiction breakdown}
\label{sec:results:perjurisdiction}

Table~\ref{tab:perjurisdiction} reports Fabrication Rate and Citation
Support F1 broken down by court tier. Outputs of the per-jurisdiction
analysis live in \texttt{citecheck/experiments/breakdowns/2027-04/}.

\begin{table}[t]
\centering
\small
\begin{tabular}{lcc}
\toprule
\textbf{Court tier} & \textbf{Fabrication $\downarrow$}
                    & \textbf{Support F1 $\uparrow$} \\
\midrule
US Supreme Court             & <TBD> & <TBD> \\
Federal Circuits (1--11)     & <TBD> & <TBD> \\
DC Circuit                   & <TBD> & <TBD> \\
Federal Circuit              & <TBD> & <TBD> \\
State appellate (if in v0.1) & <TBD> & <TBD> \\
\bottomrule
\end{tabular}
\caption{CiteCheck performance by court tier on CiteCheck v0.1.}
\label{tab:perjurisdiction}
\end{table}

\subsection{Ablations}
\label{sec:results:ablations}

Table~\ref{tab:ablations} isolates the contribution of each pipeline
component. The relevance-only ($\lambda = 0$) row isolates the
contribution of the groundedness term in Equation~\ref{eq:loss}.

\begin{table*}[t]
\centering
\small
\begin{tabular}{lccp{6cm}}
\toprule
\textbf{Configuration} &
  \textbf{Fabrication $\downarrow$} &
  \textbf{Support F1 $\uparrow$} &
  \textbf{Note} \\
\midrule
Full CiteCheck                                & <TBD> & <TBD>
   & All three stages \\
\quad $-$ Stage 2 reranker                    & <TBD> & <TBD>
   & Hybrid retrieval $\rightarrow$ Stage 3 directly \\
\quad $-$ Stage 3 \texttt{CitationResolver}   & <TBD> & <TBD>
   & Reranker + post-hoc citation check \\
\quad $-$ Stage 3 constrained decoding        & <TBD> & <TBD>
   & Free-form citation emission \\
\quad $-$ Both Stage 3 components             & <TBD> & <TBD>
   & Retrieval + reranker only;
     isolates Stage-2 contribution \\
Reranker $\lambda = 0$ (relevance only)       & <TBD> & <TBD>
   & Isolates the groundedness term in Eq.~\ref{eq:loss} \\
\bottomrule
\end{tabular}
\caption{Ablations on CiteCheck v0.1, single seed ($=42$).}
\label{tab:ablations}
\end{table*}

\subsection{Cost analysis}
\label{sec:results:cost}

Table~\ref{tab:cost} reports latency and token cost per query.
Headline summary placeholder: the verify loop adds $<TBD>$\% to p95
latency relative to Naive RAG, at a fabrication-rate reduction of
$<TBD>$ percentage points.

\begin{table}[t]
\centering
\small
\begin{tabular}{lcccc}
\toprule
\textbf{System} & \textbf{p50 (s)} & \textbf{p95 (s)}
                & \textbf{Tokens/q} & \textbf{Tools/q} \\
\midrule
Naive RAG              & <TBD> & <TBD> & <TBD> & 0     \\
Self-RAG               & <TBD> & <TBD> & <TBD> & 0     \\
CRAG                   & <TBD> & <TBD> & <TBD> & <TBD> \\
\textbf{CiteCheck}     & \textbf{<TBD>} & \textbf{<TBD>}
                       & \textbf{<TBD>} & \textbf{<TBD>} \\
\bottomrule
\end{tabular}
\caption{Latency and token cost per query.}
\label{tab:cost}
\end{table}

\subsection{Human audit}
\label{sec:results:audit}

On the 200-item human-audit subset, Citation Support F1 from the NLI
judge is $<TBD>$, the human-audited variant is $<TBD>$, and
human/NLI agreement (Cohen's $\kappa$) is $<TBD>$. On the 100-item
Answer Utility subset, CiteCheck scores $<TBD> \pm <TBD>$ versus
$<TBD> \pm <TBD>$ for the best baseline on a 5-point Likert scale.
We report both NLI and human numbers throughout; per the released
design specification we do \emph{not} headline NLI-only numbers.

% =====================================================================
%  7. Discussion
% =====================================================================
\section{Discussion}
\label{sec:discussion}

\subsection{What did and did not beat the baselines}
\label{sec:discussion:results}

\todo{Write after Phase~3 numbers land. Expected narrative: the
agentic verify loop dominates on Fabrication Rate and Citation
Resolution Rate; ablations show the reranker contributing most to
Citation Support F1 and the \texttt{CitationResolver} contributing
most to Fabrication Rate. Honest contingency: if CRAG or Self-RAG
already drives Fabrication Rate near zero on this benchmark, reframe
the contribution around the Bluebook-aware reranker and per-axis
decomposition rather than headline reduction.}

\subsection{Open vs.\ closed}
\label{sec:discussion:openclosed}

\todo{Where CiteCheck matches or exceeds GPT-4o-mini with web search,
emphasize that this is an open-weight 7--8B model on a single 24~GB
GPU with no proprietary API. Where GPT-4o-mini wins, be honest about
which axis and discuss whether the gap closes with the secondary
backbone or the 13B stretch.}

\subsection{Implications for legal-tech deployment}
\label{sec:discussion:deployment}

The deployment story has two pieces. First, the
\emph{Mata}-type harm \citep{mata2023} is precisely the failure mode
that a verify loop returning ``unresolvable citation'' rather than a
confident fabrication addresses; this is the minimum bar a deployed
system must clear. Second, the per-citation verification annotations
(\texttt{verified} / \texttt{unresolvable} / \texttt{non\_supporting})
provide the audit trail whose absence
\citet{magesh2025hallucination} flagged in closed commercial tools.

\subsection{Limitations}
\label{sec:discussion:limitations}

We list six known limitations.

\paragraph{State coverage.} v0.1 restricts to federal appellate and
the US Supreme Court; state appellate is v0.2 work, gated on uneven
CourtListener coverage at the state level.

\paragraph{NLI proxy.} Legal ``stands for'' is a strictly stronger
relation than textual entailment; we report human-audited Support F1
alongside the NLI-judged variant to bound the proxy error.

\paragraph{Grammar coverage.} The v0.1 Bluebook grammar covers
${\sim}80$\% of common case-cite forms; unsupported forms (string
cites, signal-heavy parentheticals, foreign parallel cites) are
logged in Appendix~\ref{app:grammar}.

\paragraph{Latency overhead.} The verify loop adds an estimated
$<TBD>$\% to p95 latency; deployment tolerance depends on whether
the use case is interactive chat or asynchronous brief preparation.

\paragraph{Language and jurisdiction scope.} US English only;
portability to EU, UK, or Indian corpora (cf.\ EL-RAG's COLIEE and
NyayaRAG evaluations) is not demonstrated.

\paragraph{Annotator pool.} Small ($N = <TBD>$ total) with uneven
practice-area coverage: constitutional law is over-represented
relative to administrative law, and civil procedure is
over-represented relative to trusts-and-estates.

\subsection{Ethical considerations}
\label{sec:discussion:ethics}

\paragraph{Dual-use.} A fabrication benchmark could in principle be
repurposed for adversarial fine-tuning to produce more plausible
fabrications. We mitigate via the CC-BY license terms, a README
advisory, and DOI-citable methodology so that downstream uses are
traceable.

\paragraph{Seed provenance.} \citet{charlotin2025tracker} entries
derive from public court orders. We cite docket numbers without
re-publishing attorney names beyond what is already in the public
record.

\paragraph{Not legal advice.} The benchmark and the method do not
constitute legal advice. Deployment guidance is collected in
Appendix~\ref{app:repro}.

% =====================================================================
%  8. Conclusion
% =====================================================================
\section{Conclusion}
\label{sec:conclusion}

We contribute three artifacts. \emph{First}, CiteCheck v0.1, an open
${\sim}500$-question benchmark for US case-law citation correctness
scored on three orthogonal axes -- existence, support, and
jurisdictional validity. \emph{Second}, the CiteCheck pipeline: a
three-stage agentic-RAG system that combines hybrid retrieval, a
Bluebook-structure-aware faithfulness-tuned reranker, and an agent
verify loop that resolves emitted citations against CourtListener at
decode time. \emph{Third}, an empirical comparison against five open
baselines, one closed reference, and one post-hoc detector that
isolates the contribution of every pipeline stage.

We release CiteCheck v0.1 under CC-BY, the pipeline code under MIT,
all prompts, and three random seeds at
\url{https://github.com/<TBD>}, co-located with an arXiv preprint.
In future work we will (i)~expand the benchmark to state courts in
v0.2; (ii)~extend the method to EU and UK case law; and
(iii)~integrate overruling detection
\citep{zhang2025overruled} as a fourth axis, capturing
\emph{precedent-currency validity}.

% =====================================================================
%  Acknowledgments  (after \section so it isn't numbered)
% =====================================================================
\section*{Acknowledgments}
\todo{Acknowledge advisor (TBD), annotators (by name if they
consent, by role otherwise), the Free Law Project for \texttt{eyecite}
and CourtListener, and Caselaw Access Project / Harvard Law School
Library.}

% =====================================================================
%  References
% =====================================================================
\bibliography{references}
\bibliographystyle{acl_natbib}

% =====================================================================
%  Appendices
% =====================================================================
\appendix

\section{Bluebook Grammar}
\label{app:grammar}

\todo{Reproduce the EBNF/PEG of the v0.1 Bluebook grammar from
\texttt{citecheck/src/citecheck/agent/grammar.py}. Add a coverage
table for supported vs.\ out-of-grammar forms in the CAP development
sample, with counts for the four most common unsupported cases:
string cites, signal-heavy parentheticals, foreign parallel cites,
and short cites that omit the case name.}

\section{Prompts}
\label{app:prompts}

\todo{Verbatim prompts (with temperature and max-token settings) for
every baseline and for CiteCheck's planner, rewriter, generator, and
verifier modules. List the prompt-revision history in a footnote so
reviewers can correlate prompt versions with experiment runs.}

\section{Extended Per-Task Tables}
\label{app:tables}

\todo{Full per-task and per-jurisdiction breakdowns produced by
Task~37, beyond the summary in Table~\ref{tab:perjurisdiction}.
Include per-axis CIs and the full $\lambda$-sweep curve referenced in
Section~\ref{sec:method:reranker}.}

\section{Reproducibility Checklist}
\label{app:repro}

\todo{ACL/NeurIPS-standard reproducibility checklist: compute,
hyperparameters, splits, package and model versions, seeds, the exact
\texttt{citecheck run \ldots} CLI invocations for each table, and
expected runtimes. Include the documented EL-RAG reproduction gap
and the IRCoT fallback decision rule from
Section~\ref{sec:setup:baselines}.}

\section{Annotator Instructions}
\label{app:annotation}

\todo{Full guideline text, the 50-item gold calibration set, the
adjudication protocol used when the two annotators disagreed on the
double-annotated subset, and the demographic / training summary for
the annotator pool (subject to the consent terms recorded in
Section~\ref{sec:discussion:ethics}).}

\end{document}
